\newpage
% \color{white}
% \pagecolor{black}


\ifdefined\useChapters
\chapter{O passado a quem pertence?}
 
\else
\chapter{}
\fi

\begin{center}
\textit{Garoto e Ítalo, vindos do futuro}
\end{center}

O cheiro de mato molhado é inconfundível. A umidade sobe do chão, entra pelo nariz, gruda na roupa. Estamos no meio da floresta.  
Mas como viemos parar aqui, se eu não estava desdobrado? Muito pelo contrário: eu estava bem acordado, inteiro dentro do corpo, e em um dos momentos mais tensos da minha vida.

Não era uma visão, uma lembrança projetada ou uma viagem apenas da mente. Meu corpo e o de Ítalo tinham atravessado o tempo juntos. Naquele passado, portanto, existiam duas versões de cada um de nós: as que viveram o massacre pela primeira vez e as que acabavam de chegar do futuro para tentar impedi-lo.

— Eu não acredito que você fez isso, você é a pessoa mais incrível do mundo! — fala Ítalo empolgado, pulando e gesticulando no meio da mata, como se estivéssemos num parque de diversões.

Ele gira, olha ao redor, pisa em folhas secas, cheira o ar.

— Você tem noção do que acabou de fazer? — continua. — Você tirou a gente de um lugar e trouxe pra outro totalmente diferente, como naqueles filmes de ficção científica!

Ele me dá alguns empurrões e socos no ombro, em sinal de entusiasmo.  
Eu, por outro lado, só consigo sentir um frio estranho na barriga.

— Eu não sabia que era capaz disso — digo. — Não enquanto estou no meu corpo.  
Fora dele, já fui algumas vezes pra outros lugares e tempos, mas nunca trouxe o corpo junto.

Ítalo para por um momento, como se juntasse peças.

— Você já foi pra outros tempos? — pergunta, mais sério. — Existe a chance de a gente não ter só mudado de lugar?

Fico em silêncio por alguns segundos.  
O barulho da floresta parece aumentar: pingos caindo das folhas, zumbido de insetos, o som distante de um pássaro.

— Puts — respondo, enfim. — Você tem toda razão.  
Como eu não tive controle, não sei pra onde… ou pra quando viemos.  
Geralmente, eu ia pra algum lugar e tempo que tivesse conexão com o que eu estava sentindo, mas dessa vez eu só queria estar seguro.

Ele me encara.

— Seguro? — repete. — E a sua cabeça te trouxe justo pra uma floresta no meio do nada. Ótimo critério.

Tento não rir.  
Apesar do medo, há algo familiar nesse lugar.  
O cheiro, a textura do solo, a forma como a luz passa entre as árvores.

— Vamos andar um pouco por aqui pra ver se entendemos o que aconteceu — digo.

— Só digo pra tomarmos cuidado — fala Ítalo, se esgueirando por entre as árvores de forma furtiva. — Sem querer, você pode ter nos colocado numa grande roubada.

Caminhamos por um bom tempo.  
Pra mim, tudo ali parecia igual: árvores por todos os lados, mato molhado, sombras repetidas.  
Mas aquela sensação de déjà vu não saía do peito.  
Era como estar dentro de um sonho que eu já tinha esquecido.

Depois de um tempo, ouvi passos apressados.

Não era o som leve de um animal.  
Era pesado, ritmado, múltiplo — como uma perseguição.

— Agacha aqui — sussurrei. — Tem alguém vindo.

Nos escondemos atrás de um tronco caído.  
Entre as frestas das folhas, vimos pessoas correndo por entre a floresta, em grupo, desviando de galhos, empurrando o mato.  
Estavam perseguindo alguém.

Ítalo arregalou os olhos.

— Olha ali… — murmurou. — Eles estão perseguindo alguém…

Ele apertou meu braço.

— Peraí.  
Isso ali… é você, cara.

Meu corpo gelou.

— Eu? Mas onde será… — comecei a perguntar, até que a ficha caiu. — Espera. Eu acho que sei quando estamos.

As cenas começaram a se encaixar na memória, como se alguém tivesse tirado o pó de um quadro antigo.

— Eu acho que voltamos pro momento em que vocês me perseguiam na floresta, depois de invadirem minha escola — digo, sentindo a garganta fechar.

Fico olhando o “eu passado” correndo, tropeçando, tentando fugir.  
É estranho demais ver a própria história de fora.

— Eu não posso deixar que eles morram de novo — falo, mais pra mim mesmo do que pra Ítalo. — Não posso deixar o massacre acontecer outra vez.

Saio correndo, seguindo a mim mesmo e aqueles que me perseguiam, tentando não perder ninguém de vista.  
Não lembro exatamente o caminho até o vilarejo, mas sei que o eu do passado vai encontrar.  
Então, não posso me afastar dele, nem deixar que ninguém nos veja.

— Espera aí, cara, tem uma coisa que eu preciso te contar — grita Ítalo, correndo atrás de mim.

— Agora não temos tempo pra conversar — respondo sem olhar pra trás. — Se você vai me ajudar, se apressa!

Como já consigo volitar com segurança, começo a seguir pelo alto, pairando entre as copas das árvores.  
É mais difícil ser notado lá de cima e faço menos barulho.  
Ítalo vem logo atrás, também volitando, mais experiente, mais leve.

Lá na frente, vejo a cena que já conhecia — só que dessa vez de um ângulo diferente.

O momento em que eu tropeço.

Na época, achei que tinha sido só azar.  
Agora, vejo com clareza: foi crucial.  
Eu caio em um buraco, rolo por entre várias árvores mortas em uma encosta que formava uma espécie de corredor natural.

Daqui de cima, a imagem é clara: os troncos secos, empilhados, conduzindo meu corpo até o ponto exato onde tudo mudaria.

Do lado de fora, continuo seguindo meu próprio caminho com o olhar.

Logo depois, vejo algo que, na outra vez, só percebi pela metade.

A cúpula ao redor do vilarejo.

É linda.  
Uma bolha de sabão gigante, brilhando em reflexos iridescentes.  
Pra qualquer outra pessoa, parece só ar tremendo.  
Pra mim, é um campo de força nítido, quase sólido.

Descubro mais tarde que Ítalo e os outros não conseguiam ver isso.  
Mas, por algum motivo, eu conseguia.

Logo que meu corpo parou de rolar, os moradores da tribo chegaram para me buscar.  
Eles não demoraram tanto quanto eu lembrava.  
Na minha memória, eu tinha ficado muito tempo largado no chão, sozinho.  
Vendo de fora, percebo que o tempo foi bem menor.  
Talvez a solidão tenha esticado a lembrança.

\begin{center}
  $\ast$~$\ast$~$\ast$

  \textit{Casa de Crispim, no presente}
\end{center}

Enquanto isso, na sala da casa de Crispim:

— O que aconteceu? Pra onde ele foi, Crispim? — pergunta o garçom, indignado, ao ver a arma caindo no chão depois do desaparecimento do Garoto e de Ítalo.

— Eu ainda não sei — responde Crispim, ofegante. — Mas vou atrás deles. Cuide do meu corpo enquanto eu vou.

— Cuidar do seu corpo… — o garçom revira os olhos. — Eu sempre acho isso estranho.  
Tá bom, tá bom, eu cuido. Vai logo, velho.

Crispim se aproxima da cadeira, a mesma em que esteve tantas vezes.  
Senta-se com cuidado, como se fosse um ritual antigo.

— E não pense que vai ficar por isso mesmo ter matado minha filha — diz, encarando o garçom. — Quando eu voltar, vou pensar numa punição à altura.

O garçom engole seco, mas tenta esconder.

Crispim fecha os olhos, respira fundo.  
Começa a se concentrar no Garoto e em Ítalo, tentando encontrar os rastros dos dois no tempo-espaço.

Perseguir alguém no tempo e no espaço exigia um domínio do desdobramento que Crispim possuía havia anos.
A sala vai sumindo pra ele.  
A floresta vai surgindo.

\begin{center}
  $\ast$~$\ast$~$\ast$

  \textit{Crispim, no passado do vilarejo}
\end{center}

Crispim desdobrado chega na floresta e vê Larissa, Pedro e Ítalo do passado andando em círculos, perseguindo o Garoto daquela época em vão.
Eles não faziam ideia de onde o Garoto daquela época havia parado.
Ficariam ali o dia todo, rodando, cheios de raiva e sem direção.

Ele, no entanto, sabe exatamente onde o Garoto daquela época está.
Consegue senti-lo.

— Olhem aqui, seus idiotas — grita Crispim, enquanto se materializa para que todos o vejam.

Larissa quase tropeça de susto.

— Crispim? O que está fazendo aqui? — pergunta ela, a respiração pesada.

— Agora eu não tenho tempo pra explicar — diz ele, impaciente. — Estão vendo aquelas árvores mortas amontoadas ali?

Ele aponta para a encosta por onde o corpo do Garoto daquela época havia rolado.

— Ele rolou por entre elas. Sigam a trilha.  
Quando chegarem lá embaixo, vão notar que tem algo errado com o ambiente.  
É um campo de força, e ele está lá dentro se escondendo.  
Esperem um sinal e comecem a atacar.

Assim que termina de falar, volita em direção ao vilarejo.  
Fora do corpo, ele tem muito mais liberdade — e uma frieza que incomoda.

\begin{center}
  $\ast$~$\ast$~$\ast$

  \textit{Garoto e Ítalo, vindos do futuro}
\end{center}

Não demora muito até que Larissa, Pedro e Ítalo do passado cheguem perto do vilarejo.

Dessa vez, ao contrário do que eu lembrava, eles não entram logo de cara.  
Ficam parados do lado de fora, observando, desconfiados, como se sentissem um “não” silencioso no ar.

Lá dentro, eu atravessei a cúpula de energia com facilidade.  
Aparentemente, a anciã me deixou entrar.  
Ou me reconheceu de algum jeito.

Fui correndo ao encontro dela, o coração disparado, tentando juntar em uma frase tudo o que precisava dizer.

— O que você está fazendo aqui, menino? Não era pra você estar aqui — diz a anciã assim que me aproximo.

— Como assim não era pra eu estar aqui? — pergunto, sem fôlego. — Você já me conhece?

Ela me observa de cima a baixo, com aquele olhar que parece atravessar tempo e corpo.

— Eu vi seu corpo lá com a menina — diz, com calma. — E agora você aparece aqui desse jeito, com a energia um pouco diferente… um pouco mais velha.  
O que veio fazer aqui?  
Você veio do futuro, filho?

Engulo seco.

— Vim. Como você sabe?

Ela sorri de leve.

— Sempre me esqueço que você sabe de muita coisa — completo, meio sem graça.

Ela se aproxima, pega minhas mãos nas dela.

— Sabe, filho… a gente pode mudar o futuro com as nossas ações, mas nunca pode mudar o passado. Ele já foi escrito.  
Eu não quero saber o que vai acontecer aqui pra você ter voltado. Não me conte.  
Mas tome cuidado com o que vai fazer.

— Se esconde — peço, quase implorando. — Eu vou proteger vocês.

Do lado de fora, vejo Ítalo balançando os braços, olhando pra mim de forma frenética, como se dissesse “não faz isso”.  
Eu ignoro.  
Ou talvez escolha não entender.

Logo, os três — Larissa, Pedro e o Ítalo do passado — começam a se aproximar da cúpula.  
Antes que eles nos vejam, eu ataco.

Labaredas de fogo saem das minhas mãos e atingem a barreira.  
Larissa, lá de fora, revida com chamas ainda mais fortes.

Só então me dou conta:  
Até eu atacar, eles não estavam nos vendo.  
Eu mesmo tinha nos denunciado.

Como eu sou burro!

Eles começam o ataque com força total.  
As rajadas de energia racham a cúpula.  
Cada fissura é como um estalo no meu peito.

Não demora para que encontrem um ponto fraco e consigam romper.

A barreira cai.  
Eles entram.

Nesse momento, Ítalo do futuro também atravessa a cúpula, finalmente conseguindo entrar.  
Lutamos juntos contra os três.

Dessa vez, eu estava um pouco mais preparado.  
Sabia o que estava por vir.  
Tinha treinado com Ítalo.  
Tinha visto aquele futuro uma vez.

Mas não foi suficiente.

No meio da batalha, todos os homens e mulheres do vilarejo se juntam a nós.  
Lutam com o que têm: paus, ferramentas, coragem.  
Têm mais coração do que poder.

Ainda assim, o massacre acontece.  
De novo.

Foi tudo muito rápido.  
Mesmo sabendo o que aconteceria, não fomos capazes de mudar o destino.

Quando percebi, eu e o Ítalo que viera comigo do futuro estávamos desacordados.
Caídos no chão, como na primeira vez.

Acordei com a cabeça pesada, o corpo doendo, a visão embaçada.

A primeira coisa que vi foi a anciã segurando a menina no colo.  
Ela estava triste, mas com um semblante calmo — como quem vê uma cena repetida e, ainda assim, escolhe ficar.

— Me desculpe — digo, com dificuldade. — Eu não pude mudar nada, como imaginava.

— Tudo bem, filho — ela responde, fazendo carinho na cabeça da menina. — Como eu te disse, tem coisas que têm que ser.

Ao contrário do que eu imaginava, a menina ainda não estava morta.  
A respiração era fraca, mas existia.

— O que é isso, vovó? — ela pergunta, com a voz quase apagada, mal abrindo os olhos. — Ele não acabou de sair correndo daqui?

Ela me reconhece.  
Como se duas lembranças brigassem dentro dela.

— Você consegue me ouvir? — pergunta, olhando bem fundo nos meus olhos.

— Consigo — respondo, engolindo o choro.

Ela sorri, cansada, mas inteira.

— Olha só, minha vó sempre me ensinou que a gente tem que ser grato pelas situações da vida — diz. — E eu sou grata por esse momento de vida que ainda tenho.

As palavras dela atravessam tudo: culpa, medo, cansaço.

— Quando eu te curei — continua — eu estava na verdade te doando amor.  
Mentalizando amor pra você, direcionado pras partes que você tinha falta dele em forma de ferimentos.

Ela coloca as mãos no meu peito.  
Uma paz imensa me atravessa, como se alguém tivesse acendido luz dentro de um quarto esquecido.

Sinto o peso nas costas diminuir.  
Não porque a responsabilidade tenha sumido, mas porque, por um instante, não estou sozinho dentro dela.

— Você se lembra quando eu te perguntei se você se lembrava do que tinha acontecido nesse dia? — pergunta Ítalo, se aproximando.

— Da outra vez, eu me lembro de achar que tinha te visto lutar contra a gente — respondo. — Mas foi tudo tão intenso e escuro, por causa da mata fechada, que eu não tinha certeza.  
Mas hoje eu tenho.  
Isso que aconteceu agora… sempre aconteceu.  
A gente veio até aqui porque tinha que vir, de certa forma.

Ele passa a mão no cabelo, pensativo.

— Não temos tanto poder de mudar o tempo como eu imaginava — diz. — O que faz com que aquele futuro que você viu seja ainda mais provável de acontecer.

Ficamos em silêncio por alguns segundos.

— Mas, como a anciã disse, o futuro a gente pode moldar com os nossos atos — ele completa. — Então… ainda há uma chance.

\begin{center}
  $\ast$~$\ast$~$\ast$

  \textit{Garçom, na casa de Crispim, no presente}
\end{center}
   
Enquanto isso, o garçom estava todo esse tempo ao lado do corpo de Crispim, na sala.  
Observava o velho imóvel, recostado na cadeira, com uma expressão neutra que escondia tudo.

Até que vislumbrou uma oportunidade.

Aproximou-se devagar.

— Você acha mesmo que pode me ameaçar e sair impune? — murmurou, sabendo que Crispim não podia ouvir.

Estendeu a mão sobre o peito dele.

Usou seu poder de retirar as habilidades dos outros.  
Dessa vez, em Crispim.

Em poucos instantes, o elo entre corpo e espírito se rompeu.  
Crispim não conseguiu mais voltar.  
Ficou preso do lado de fora, vagando.  
Sem corpo.  
Sem âncora.

O garçom sorriu, aliviado.

Tirava de cena aquele que já não via mais como líder.  
Um trono vazio começava a se formar — e ele sabia exatamente quem pretendia sentar nele.
